\section{I/O}

\begin{question}
	以下哪个类用于实现文件读取？
	\begin{options}
		\item FileReader
		\item FileWriter
		\item BufferedReader
		\item BufferedWriter
	\end{options}
	\begin{answer}
		A
	\end{answer}
	\begin{explanation}
		\texttt{FileReader} 用于读取字符文件。
	\end{explanation}
\end{question}

\begin{question}
	以下哪个类用于实现文件写入？
	\begin{options}
		\item FileReader
		\item FileWriter
		\item BufferedReader
		\item BufferedWriter
	\end{options}
	\begin{answer}
		B
	\end{answer}
	\begin{explanation}
		\texttt{FileWriter} 用于写入字符文件。
	\end{explanation}
\end{question}

\begin{question}
	以下关于Java中的流的说法错误的是
	\begin{options}
		\item 分为字节流和字符流
		\item 字节流以字节为单位处理数据
		\item 字符流以字符为单位处理数据
		\item 字节流和字符流可以相互转换
	\end{options}
	\begin{answer}
		D
	\end{answer}
	\begin{explanation}
		字节流和字符流不能直接相互转换，需通过编码转换（如 \texttt{InputStreamReader}）。
	\end{explanation}
\end{question}

\begin{question}
	以下哪个方法用于创建一个新文件？
	\begin{options}
		\item createNewFile()
		\item mkdir()
		\item mkdirs()
		\item exists()
	\end{options}
	\begin{answer}
		A
	\end{answer}
	\begin{explanation}
		\texttt{File.createNewFile()} 创建新文件。
	\end{explanation}
\end{question}

\begin{question}
	以下哪个方法用于删除一个文件？
	\begin{options}
		\item delete()
		\item remove()
		\item clear()
		\item erase()
	\end{options}
	\begin{answer}
		A
	\end{answer}
	\begin{explanation}
		\texttt{File.delete()} 删除文件或目录。
	\end{explanation}
\end{question}

\begin{question}
	以下关于Java中序列化的说法错误的是
	\begin{options}
		\item 可以将对象转换为字节序列
		\item 可以通过实现Serializable接口实现序列化
		\item 静态成员变量不能被序列化
		\item 所有成员变量都会被序列化
	\end{options}
	\begin{answer}
		D
	\end{answer}
	\begin{explanation}
		被 \texttt{transient} 修饰的变量不会被序列化。
	\end{explanation}
\end{question}

\begin{question}
	以下哪个方法用于反序列化？
	\begin{options}
		\item readObject()
		\item deserialize()
		\item restoreObject()
		\item recoverObject()
	\end{options}
	\begin{answer}
		A
	\end{answer}
	\begin{explanation}
		\texttt{ObjectInputStream.readObject()} 用于反序列化。
	\end{explanation}
\end{question}
